% Seccion 8 modulo de recuperacion de texto

\section{Módulo de Recuperación de Texto: Índice Invertido en Disco}

Una búsqueda de texto funciona distinto a las consultas del primer entregable.
No compara una columna contra un valor, sino que pide los documentos que
contienen ciertas palabras y, además, quiere saber cuáles son los más
relevantes. Para responder esto sin leer la tabla completa construimos un índice
invertido que se guarda en disco con las mismas reglas que el resto del motor:
bloques de 4\,096 bytes, acceso con \texttt{seek} y \texttt{struct}, y cada
bloque leído o escrito contado por el \texttt{DiskCounter}.

\subsection{Preprocesamiento del Texto}

Los documentos y las consultas siempre pasan por el mismo analizador. Si no
fuera así, una palabra de la consulta podría no coincidir nunca con la forma en
que se guardó en el índice. El analizador aplica cinco pasos:

\begin{enumerate}
    \item \textbf{Minúsculas (case folding):} todo el texto pasa a minúsculas.
    \item \textbf{Tokenización:} el texto se separa en palabras; la puntuación
    sirve de separador y los números sueltos se descartan.
    \item \textbf{Normalización:} se quitan las tildes.
    \item \textbf{Palabras vacías (stopwords):} se descartan las 186 palabras de
    la lista en español (\texttt{stoplist\_es.txt}), como "el", "de" o "que".
    \item \textbf{Stemming:} cada palabra se reduce a su raíz con el algoritmo
    Snowball en español de \texttt{nltk}. Por ejemplo, "cámaras" queda en
    \texttt{cam} y "verificación" en \texttt{verif}.
\end{enumerate}

El orden entre el tercer y el quinto paso lo descubrimos probando. Al principio
sacábamos la raíz de la palabra con tilde y quitábamos la tilde después. Con eso
"fotografía" quedaba en \texttt{fotograf}, pero "fotografia", escrita sin tilde
en una consulta, quedaba en \texttt{fotografi}, así que las dos nunca
coincidían. Al quitar las tildes antes de sacar la raíz, el resultado ya no
depende de cómo escriba el usuario.

\subsection{Diseño Físico del Índice}

El índice se guarda en tres archivos binarios. La página 0 de cada uno es una
cabecera con un identificador del archivo y los totales; el resto de páginas
guarda entradas de tamaño fijo:

\begin{table}[H]
\centering
\caption{Archivos del índice invertido}
\begin{tabularx}{\textwidth}{@{}llX@{}}
\toprule
\textbf{Archivo} & \textbf{Entrada} & \textbf{Contenido} \\
\midrule
\texttt{vocab.dic} & 52 B & \texttt{<término 32s, df, cf, offset, nbytes>}, ordenadas por
término, 78 por página. Al final, un directorio con el primer término de cada página \\
\texttt{postings.inv} & 8 B & listas de posteo \texttt{<doc\_id, tf>} una tras otra,
cada una ordenada por \texttt{doc\_id} \\
\texttt{docs.bin} & 20 B & \texttt{<RID, norma L2, longitud>} de cada documento,
204 por página \\
\bottomrule
\end{tabularx}
\end{table}

Tres decisiones determinan cuánto cuesta una consulta:

\begin{itemize}
    \item \textbf{Encontrar un término cuesta una sola lectura.} Al abrir el
    índice se carga en memoria el directorio de \texttt{vocab.dic}, igual que se
    hace con la raíz de un árbol B$^+$. Con ese directorio se sabe en qué página
    está el término, y dentro de esa página se busca con búsqueda binaria.
    \item \textbf{Las listas se guardan una detrás de otra, sin dejar espacio
    vacío.} Un término que aparece en tres documentos ocupa 24 bytes, no una
    página entera. Por eso leer una lista cuesta solo las páginas que ocupa.
    \item \textbf{La norma de cada documento se calcula al construir el índice.}
    Para calcular el coseno se necesita la norma de cada documento candidato. Al
    tenerla guardada en \texttt{docs.bin} se puede ordenar por relevancia sin
    leer la tabla, y esto es lo que más marca la diferencia frente a PostgreSQL
    en el experimento 6.
\end{itemize}

El índice guarda los RID de las filas de la tabla, así que después de buscar en
el índice cada resultado se lee de la tabla por su RID, igual que en cualquier
índice secundario.

\subsection{Construcción con SPIMI}

El índice se construye recorriendo la colección una sola vez, con el método
SPIMI (\textit{Single-Pass In-Memory Indexing}):

\begin{enumerate}
    \item Cada documento pasa por el analizador y sus términos se agregan a un
    diccionario en memoria, donde cada término guarda su lista
    \texttt{[(doc\_id, tf)]}.
    \item Cuando el diccionario llega a 200\,000 posteos, se ordena por término y
    se guarda en disco como un bloque, también en páginas de 4\,096 bytes.
    \item Al terminar, se juntan todos los bloques en uno solo usando un heap
    (\textit{k-way merge}). Por cada término se unen sus posteos, se escriben en
    \texttt{postings.inv}, se agrega su entrada a \texttt{vocab.dic} y se suma su
    parte a la norma de cada documento.
    \item Los tres archivos nuevos reemplazan a los anteriores de una sola vez,
    para que el índice nunca quede a medio escribir.
\end{enumerate}

Agregar documentos a un índice que ya existe usa el mismo proceso: los bloques
nuevos se juntan con el índice actual, que ya está ordenado por término y se
trata como un bloque más. Así no hace falta volver a analizar los documentos que
ya estaban indexados.

\subsection{Consultas Booleanas y Ranking por Coseno}

El peso de un término $t$ en un documento $d$ es
\begin{equation}
    w_{t,d} = \left(1 + \log_{10} \text{tf}_{t,d}\right) \cdot \log_{10} \frac{N}{\text{df}_t},
\end{equation}
y la relevancia de $d$ para la consulta $q$ es la similitud de coseno
\begin{equation}
    \cos(q, d) = \frac{\sum_{t \in q} w_{t,q}\, w_{t,d}}{\lVert q \rVert_2 \, \lVert d \rVert_2}.
\end{equation}
Solo se recorren las listas de los términos que aparecen en la consulta, y los
$K$ documentos con mayor puntaje se obtienen con un heap.

Las consultas booleanas aceptan \texttt{AND}, \texttt{OR}, \texttt{NOT} y
paréntesis, y \texttt{NOT} se evalúa primero, luego \texttt{AND} y al final
\texttt{OR}. Cada operador recorre dos listas ordenadas al mismo tiempo con dos
punteros, así que su costo es $O(n + m)$: \texttt{AND} se queda con los
documentos que están en ambas listas, \texttt{OR} con los que están en alguna y
\texttt{AND NOT} con los de la primera que no están en la segunda. Si la
consulta no tiene operadores, se trata como texto libre: basta con que un
documento tenga alguna de las palabras, y el coseno decide el orden.

Para comprobar que el índice funciona bien, programamos también la misma
búsqueda sin índice, que recorre todos los documentos y los analiza uno por
uno. Las pruebas confirman que los dos caminos devuelven los mismos documentos
con los mismos puntajes en diez consultas de todos los tipos.

\subsection{Integración con el Motor}

\begin{lstlisting}[language=SQL, caption={Índice invertido y consulta textual}]
CREATE INDEX ix_cuerpo ON expedientes (cuerpo)
  USING INVERTED WITH (scoring = 'TFIDF_COSINE', language = 'spanish');

SELECT id, tipo, COSINE_SCORE() AS relevancia
FROM expedientes
WHERE cuerpo MATCH 'verificacion AND rostro AND NOT mascarilla'
ORDER BY relevancia DESC
LIMIT 10;
\end{lstlisting}

\begin{table}[H]
\centering
\caption{Plan elegido para una consulta textual}
\begin{tabularx}{\textwidth}{@{}lX@{}}
\toprule
\textbf{Situación} & \textbf{Plan} \\
\midrule
La columna tiene un índice \texttt{INVERTED} & \texttt{InvertedIndexScan}, y luego una lectura por cada RID \\
La columna no tiene índice & \texttt{FullTextScan}: lee todas las páginas y analiza cada texto \\
\texttt{MATCH} junto a otra condición & se busca primero con el \texttt{MATCH} y la otra condición se aplica después (\texttt{Filter}) \\
Orden por el puntaje con \texttt{LIMIT} $k$ & el índice devuelve directamente los $k$ mejores (Top-$K$) \\
\bottomrule
\end{tabularx}
\end{table}

El plan muestra por separado los bloques leídos del índice y los de la tabla, y
los del índice se dividen en vocabulario, listas de posteo y normas. Las filas
que se insertan después de crear el índice se agregan como un bloque nuevo antes
de la siguiente búsqueda. Cuando se borra una fila, su entrada se queda en el
índice, pero al leerla por su RID la tabla no devuelve nada y se descarta.

En la primera versión, después de leer cada fila volvíamos a analizar su texto
para confirmar que cumplía la consulta. Lo hacíamos por si el índice apuntaba a
una fila borrada, pero esas filas ya se descartan al leerlas, y esa revisión
extra era lo que más tiempo tomaba en las consultas que devuelven muchos
documentos. Por eso la quitamos.

\subsection{Corpus de Expedientes}

Para tener texto en el cual buscar, generamos 25\,000 expedientes en español
asociados a las identidades del sistema biométrico. Usamos una semilla fija para
que siempre se genere el mismo corpus. Hay ocho tipos de expediente: registros
de acceso, verificaciones biométricas, renovaciones de credencial, fallas
técnicas de cámaras, visitas, auditorías, solicitudes de protección de datos y
alertas. Cada persona aparece con un código (\texttt{P00042}), sin nombres. En
promedio cada expediente ocupa 470 bytes.

La primera versión del corpus tenía solo 279 palabras distintas, y ese número no
aumentaba al agregar más documentos. Eso no se parece a un texto real y habría
dado resultados engañosos en el experimento. Por eso agregamos observaciones con
un vocabulario más amplio, donde pocas palabras se repiten mucho y la mayoría
aparece poco (distribución de tipo Zipf), además de códigos únicos de equipos y
tickets. Con eso el vocabulario sigue creciendo al agregar documentos, aunque
cada vez más despacio, como describe la ley de Heaps: pasa de 1\,791 términos
con mil documentos a 21\,412 con veinticinco mil.

\subsection{Experimento 6: Índice Invertido frente a Full Text Scan y GIN}

\begin{lstlisting}[language=bash, caption={Ejecución del experimento}]
python benchmarks/experimento6_textual.py --pg-dsn postgresql://...
\end{lstlisting}

El script construye el índice con colecciones cada vez más grandes y luego
ejecuta 60 consultas: \texttt{AND}, \texttt{OR}, \texttt{AND NOT} y Top-10 por
relevancia, cada una con 1, 3 y 5 términos y cinco consultas distintas por
combinación. Cada consulta se ejecuta por SQL de tres formas:

\begin{itemize}
    \item sobre la tabla \texttt{expedientes} del motor, sin índice
    (\texttt{FullTextScan});
    \item sobre la misma tabla, con el índice invertido;
    \item en PostgreSQL 18, con \texttt{to\_tsvector('spanish', cuerpo)} y un
    índice GIN.
\end{itemize}

La tabla del motor tiene 20\,000 filas en 6\,666 páginas de 4\,096 bytes. En
PostgreSQL los bloques son de 8\,192 bytes y los valores se obtuvieron con
\texttt{EXPLAIN (ANALYZE, BUFFERS)}, con los datos ya cargados en caché. Cada
valor es el promedio de las cinco consultas de su combinación.

\subsubsection{Construcción y tamaño en disco}

\begin{table}[H]
\centering
\caption{Construcción del índice invertido frente al índice GIN.
\textbf{Bloques} son los bloques intermedios que escribió SPIMI antes de
juntarlos.}
\begin{tabular}{@{}rrrrrrr@{}}
\toprule
\textbf{N} & \textbf{Términos} & \textbf{Bloques} & \textbf{SPIMI (s)} &
\textbf{Índice (KiB)} & \textbf{GIN (s)} & \textbf{GIN (KiB)} \\
\midrule
1\,000 & 1\,791 & 1 & 0{,}37 & 352 & 0{,}12 & 224 \\
2\,000 & 3\,033 & 1 & 0{,}53 & 660 & 0{,}12 & 392 \\
5\,000 & 6\,204 & 1 & 1{,}31 & 1\,556 & 0{,}31 & 880 \\
10\,000 & 10\,719 & 2 & 2{,}44 & 3\,024 & 0{,}66 & 1\,552 \\
20\,000 & 18\,082 & 4 & 4{,}90 & 5\,844 & 1{,}41 & 3\,096 \\
25\,000 & 21\,412 & 4 & 6{,}01 & 7\,240 & 2{,}40 & 3\,888 \\
\bottomrule
\end{tabular}
\end{table}

La columna del índice suma \texttt{vocab.dic} y \texttt{postings.inv}.
\texttt{docs.bin} ocupa 496~KiB adicionales con 25\,000 documentos.

\begin{figure}[H]
\centering
\begin{tikzpicture}
\begin{groupplot}[figura,
    group style={group size=2 by 1, horizontal sep=1.6cm},
    width=0.5\linewidth, height=6cm,
    xlabel={Documentos ($N$)},
    xtick={1000,10000,20000,25000},
    xticklabels={1\,000,10\,000,20\,000,25\,000},
    xticklabel style={rotate=30, anchor=east},
    scaled x ticks=false,
    legend pos=north west,
]
\nextgroupplot[ylabel={Tiempo de construcción (s)}]
\addplot[cSeq, mark=triangle*] coordinates {(1000,0.37)(2000,0.53)(5000,1.31)(10000,2.44)(20000,4.90)(25000,6.01)};
\addlegendentry{SPIMI}
\addplot[cHash, mark=diamond*] coordinates {(1000,0.12)(2000,0.12)(5000,0.31)(10000,0.66)(20000,1.41)(25000,2.40)};
\addlegendentry{GIN}
\nextgroupplot[ylabel={Tamaño en disco (MiB)}]
\addplot[cSeq, mark=triangle*] coordinates {(1000,0.34)(2000,0.64)(5000,1.52)(10000,2.95)(20000,5.71)(25000,7.07)};
\addlegendentry{vocab + postings}
\addplot[cHash, mark=diamond*] coordinates {(1000,0.22)(2000,0.38)(5000,0.86)(10000,1.52)(20000,3.02)(25000,3.80)};
\addlegendentry{GIN}
\end{groupplot}
\end{tikzpicture}
\caption{Tiempo de construcción y tamaño en disco según el número de
documentos. Los dos índices crecen de forma lineal, pero GIN lo hace más
despacio.}
\label{fig:e6-construccion}
\end{figure}

Los dos índices crecen de forma lineal con el número de documentos. El índice
invertido ocupa unos 296 bytes por documento, y el 85 por ciento de ese espacio
es \texttt{postings.inv}, porque cada posteo ocupa ocho bytes sin comprimir.
\texttt{vocab.dic} crece más lento porque el número de palabras distintas
también crece más lento que el número de documentos.

GIN ocupa más o menos la mitad y se construye entre 2,5 y 4 veces más rápido.
La diferencia no está en el método sino en cómo se guardan los datos: GIN
comprime sus listas (codificación delta) y está programado en C, mientras que
nuestras listas guardan los números enteros de 32 bits sin comprimir. Comprimir
las listas de la misma forma es la mejora más directa que queda pendiente, y
reduciría tanto el tamaño como las lecturas.

\subsubsection{Consultas}

\begin{table}[H]
\centering
\caption{Bloques leídos y tiempo promedio con $N = 20\,000$. En el índice
invertido, los bloques se separan en índice + tabla.}
\begin{tabular}{@{}lrrrrrrr@{}}
\toprule
& & \multicolumn{2}{c}{\textbf{Full Text Scan}} &
\multicolumn{2}{c}{\textbf{Índice invertido}} &
\multicolumn{2}{c}{\textbf{GIN}} \\
\cmidrule(lr){3-4}\cmidrule(lr){5-6}\cmidrule(l){7-8}
\textbf{Consulta} & \textbf{Filas} & \textbf{Bloques} & \textbf{ms} &
\textbf{Bloques} & \textbf{ms} & \textbf{Bloques} & \textbf{ms} \\
\midrule
AND, 1 & 5\,900 & 6\,666 & 8\,037 & 112 + 5\,899 & 4\,515 & 2\,055 & 3{,}7 \\
AND, 3 & 25 & 6\,666 & 5\,280 & 45 + 25 & 10{,}6 & 35 & 0{,}2 \\
AND, 5 & 1 & 6\,666 & 5\,294 & 48 + 1 & 9{,}9 & 19 & 0{,}4 \\
OR, 1 & 2\,484 & 6\,666 & 12\,088 & 54 + 2\,484 & 1\,801 & 698 & 2{,}7 \\
OR, 3 & 8\,070 & 6\,666 & 5\,390 & 121 + 8\,070 & 2\,000 & 2\,320 & 5{,}9 \\
OR, 5 & 15\,429 & 6\,666 & 9\,564 & 155 + 15\,428 & 3\,386 & 2\,688 & 20{,}9 \\
AND NOT, 1 & 16\,709 & 6\,666 & 8\,265 & 107 + 16\,707 & 9\,386 & 2\,688 & 11{,}8 \\
AND NOT, 3 & 184 & 6\,666 & 5\,198 & 85 + 184 & 41{,}7 & 169 & 0{,}2 \\
AND NOT, 5 & 0 & 6\,666 & 11\,329 & 45 + 0 & 6{,}2 & 17 & 0{,}1 \\
Top-10, 1 & 10 & 6\,666 & 8\,835 & 113 + 10 & 14{,}2 & 1\,933 & 8{,}1 \\
Top-10, 3 & 10 & 6\,666 & 5\,456 & 129 + 10 & 25{,}1 & 2\,345 & 20{,}6 \\
Top-10, 5 & 10 & 6\,666 & 16\,053 & 153 + 10 & 34{,}7 & 2\,688 & 37{,}4 \\
\bottomrule
\end{tabular}
\end{table}

\begin{figure}[H]
\centering
\begin{tikzpicture}
\begin{axis}[figura,
    height=7cm,
    ybar, bar width=3.2pt,
    ymode=log, log origin=infty,
    ymin=10, ymax=40000,
    ylabel={Bloques leídos},
    symbolic x coords={A1,A3,A5,O1,O3,O5,N1,N3,N5,T1,T3,T5},
    xtick=data,
    xticklabels={AND 1,AND 3,AND 5,OR 1,OR 3,OR 5,NOT 1,NOT 3,NOT 5,Top 1,Top 3,Top 5},
    xticklabel style={rotate=40, anchor=east},
    ytick={10,100,1000,10000},
    yticklabels={10,100,$10^3$,$10^4$},
    enlarge x limits=0.05,
    legend columns=3,
    legend style={at={(0.5,1.03)}, anchor=south},
]
\addplot[cHeap, fill=cHeap!70] coordinates {(A1,6666)(A3,6666)(A5,6666)(O1,6666)(O3,6666)(O5,6666)(N1,6666)(N3,6666)(N5,6666)(T1,6666)(T3,6666)(T5,6666)};
\addlegendentry{Full Text Scan}
\addplot[cSeq, fill=cSeq!70] coordinates {(A1,6011)(A3,70)(A5,49)(O1,2538)(O3,8191)(O5,15583)(N1,16815)(N3,269)(N5,45)(T1,123)(T3,139)(T5,163)};
\addlegendentry{Índice invertido}
\addplot[cHash, fill=cHash!70] coordinates {(A1,2055)(A3,35)(A5,19)(O1,698)(O3,2320)(O5,2688)(N1,2688)(N3,169)(N5,17)(T1,1933)(T3,2345)(T5,2688)};
\addlegendentry{GIN (bloques de 8 KB)}
\end{axis}
\end{tikzpicture}
\caption{Bloques leídos por tipo de consulta y número de términos, en escala
logarítmica. NOT corresponde a las consultas \texttt{AND NOT}. Sin índice
siempre se leen las 6\,666 páginas; con el índice invertido se leen hasta cien
veces menos cuando la consulta devuelve pocos documentos, pero más cuando
devuelve una buena parte de la tabla.}
\label{fig:e6-bloques}
\end{figure}

\begin{figure}[H]
\centering
\begin{tikzpicture}
\begin{axis}[figura,
    height=7cm,
    ybar, bar width=3.2pt,
    ymode=log, log origin=infty,
    ymin=0.05, ymax=60000,
    ylabel={Tiempo promedio (ms)},
    symbolic x coords={A1,A3,A5,O1,O3,O5,N1,N3,N5,T1,T3,T5},
    xtick=data,
    xticklabels={AND 1,AND 3,AND 5,OR 1,OR 3,OR 5,NOT 1,NOT 3,NOT 5,Top 1,Top 3,Top 5},
    xticklabel style={rotate=40, anchor=east},
    ytick={0.1,1,10,100,1000,10000},
    yticklabels={0{,}1,1,10,100,$10^3$,$10^4$},
    enlarge x limits=0.05,
    legend columns=3,
    legend style={at={(0.5,1.03)}, anchor=south},
]
\addplot[cHeap, fill=cHeap!70] coordinates {(A1,8037)(A3,5280)(A5,5294)(O1,12088)(O3,5390)(O5,9564)(N1,8265)(N3,5198)(N5,11329)(T1,8835)(T3,5456)(T5,16053)};
\addlegendentry{Full Text Scan}
\addplot[cSeq, fill=cSeq!70] coordinates {(A1,4515)(A3,10.6)(A5,9.9)(O1,1801)(O3,2000)(O5,3386)(N1,9386)(N3,41.7)(N5,6.2)(T1,14.2)(T3,25.1)(T5,34.7)};
\addlegendentry{Índice invertido}
\addplot[cHash, fill=cHash!70] coordinates {(A1,3.7)(A3,0.2)(A5,0.4)(O1,2.7)(O3,5.9)(O5,20.9)(N1,11.8)(N3,0.2)(N5,0.1)(T1,8.1)(T3,20.6)(T5,37.4)};
\addlegendentry{GIN}
\end{axis}
\end{tikzpicture}
\caption{Tiempo promedio por tipo de consulta, en escala logarítmica. En el
Top-10 de tres y cinco términos el índice invertido y GIN tardan casi lo mismo.}
\label{fig:e6-latencia}
\end{figure}

Sin índice, todas las consultas cuestan lo mismo: el \texttt{FullTextScan} lee
las 6\,666 páginas y analiza los 20\,000 textos sin importar cuántos documentos
devuelva, y por eso tarda varios segundos. Los tiempos del motor cambiaron un
poco entre una ejecución y otra según qué tan ocupada estaba la computadora,
pero las lecturas siempre fueron las mismas, porque las cuenta el
\texttt{DiskCounter}.

Cuando la consulta devuelve pocos documentos, el índice invertido es mucho más
rápido. En \texttt{AND} y \texttt{AND NOT} de tres y cinco términos lee entre 45
y 269 páginas en vez de 6\,666, y responde en milisegundos en lugar de segundos.
En el Top-10 solo se leen de la tabla las diez filas que se devuelven. El resto
son lecturas del índice (las listas de posteo y las normas de los candidatos),
por eso el costo sube un poco con cada término que se agrega: 123, 139 y 163
bloques.

Cuando la consulta devuelve una buena parte de la tabla, el índice resulta más
lento que no usarlo. Pasa con \texttt{OR} de palabras muy comunes y con
\texttt{AND NOT} de un solo término, que equivale a pedir todos los documentos
que no tienen esa palabra. El índice en sí lee poco, entre 54 y 155 páginas; el
problema está en la tabla. Cada fila del resultado se lee por separado con su
RID, y 15\,428 lecturas sueltas son más que las 6\,666 páginas de leer la tabla
completa. Es lo mismo que pasó con el árbol B$^+$ en el experimento 3, y por la
misma razón. Por ahora el planificador usa el índice siempre que existe. Como
\texttt{vocab.dic} ya guarda en cuántos documentos aparece cada término
(\texttt{df}), el planificador podría usar ese dato para estimar el costo y
elegir leer la tabla completa en estos casos.

Frente a GIN, en las consultas booleanas gana PostgreSQL. Es más rápido en todas
y lee menos en las que devuelven muchos documentos, porque tiene un buffer pool,
comprime sus listas y usa un \textit{bitmap heap scan}, que lee cada página de la
tabla una sola vez en lugar de una vez por fila. En el ranking el resultado
cambia. En el Top-10 de tres y cinco términos nuestro índice tarda casi lo mismo
que GIN (25 frente a 21~ms y 35 frente a 37~ms), pero lee entre 15 y 17 veces
menos bloques: 139 frente a 2\,345 y 163 frente a 2\,688. Esto pasa porque GIN
solo encuentra los documentos; para ordenarlos, \texttt{ts\_rank} tiene que leer
de la tabla el \texttt{tsvector} de cada candidato y ordenarlos todos antes de
aplicar el \texttt{LIMIT}. Nuestro índice, en cambio, ordena sin leer la tabla
gracias a las normas guardadas en \texttt{docs.bin}, y esa fue la decisión de
diseño que más ayudó.

En algunas consultas el número de filas del motor y el de PostgreSQL no coincide
exactamente, por ejemplo 8\,070 frente a 7\,984 en \texttt{OR} de tres términos.
Esto se debe a que las dos configuraciones en español no son iguales: nuestro
analizador quita las tildes antes de sacar la raíz y usa su propia lista de
palabras vacías.